\chapter{Results}
\label{chap:results}

This chapter presents the empirical evaluation of the implemented ZGC weak-reference optimisations and the exploratory weak-field mechanism described in \cref{chap:method}.

\section{Compared Configurations}
\label{sec:results-configurations}

For \texttt{WeakReference}-based processing, the comparison covers the baseline (\texttt{none}) and all seven non-empty combinations of the three optimisation flags (clear path, separate lists, dynamic array). In addition, a dedicated \texttt{weak\_fields} configuration is included, where weak-field benchmark classes are executed in the same harness and analysed with the same tooling.

This configuration set allows two complementary comparisons:
\begin{itemize}
	\item \textbf{Intra-model comparison} among \texttt{WeakReference} variants to isolate the impact of each optimisation and their interactions.
	\item \textbf{Inter-model comparison} between optimised \texttt{WeakReference} processing and weak-field representation to evaluate structural overhead trade-offs.
\end{itemize}

\section{Measured Metrics}
\label{sec:results-metrics}

The primary performance metrics are the GC timing statistics extracted from ZGC logs, including major collection time, old-generation collection time, concurrent mark-follow subphases, and concurrent reference-processing subphases. Memory behaviour is characterised using continuous samples of RSS, GC committed memory, and Java heap usage.

Following the methodology in \cref{subsec:measurement-methodology}, single-object benchmark comparisons use the per-run maximum values as primary statistics, while multi-object comparisons use per-run averages. This choice reflects the fact that single-object runs are designed to create concentrated peak load in one collection event, whereas multi-object runs distribute work over repeated phases.

\section{Presentation of Results}
\label{sec:results-presentation}

\Cref{fig:metric-boxplots-single,fig:metric-boxplots-multi} show box plots of the GC timing metrics for the single-object and multi-object benchmarks. They summarise the distribution of per-run maxima for the single-object benchmark and per-run averages for the multi-object benchmark, allowing comparison of both central tendency and variability. \Cref{fig:continuous-memory-single,fig:continuous-memory-multi} show representative continuous-memory traces for the same benchmark families, illustrating how memory usage evolves over time and across GC phases.

\begin{figure}
	\centering
	\includegraphics[width=\textwidth]{metric_boxplots_single_1.pdf}
	\caption[GC metric box plots (single-object)]{Box plots of GC metrics across evaluated configurations for the single-object benchmark.}
	\label{fig:metric-boxplots-single}
\end{figure}

\begin{figure}
	\centering
	\includegraphics[width=\textwidth]{metric_boxplots_double_2.pdf}
	\caption[GC metric box plots (multi-object)]{Box plots of GC metrics across evaluated configurations for the multi-object benchmark.}
	\label{fig:metric-boxplots-multi}
\end{figure}

\begin{figure}
	\centering
	\includegraphics[width=\textwidth]{continuous_memory_single_1.pdf}
	\caption[Continuous memory trace (single-object)]{Continuous memory trace for the single-object benchmark.}
	\label{fig:continuous-memory-single}
\end{figure}

\begin{figure}
	\centering
	\includegraphics[width=\textwidth]{continuous_memory_double_2.pdf}
	\caption[Continuous memory trace (multi-object)]{Continuous memory trace for the multi-object benchmark.}
	\label{fig:continuous-memory-multi}
\end{figure}
